\documentclass{article}
\usepackage{amsmath,amssymb}
\usepackage{lmodern}
\usepackage{iftex}

\author{}
\date{}

\begin{document}

\newpage
\section{Question 2}

\subsection{Part 2}

Mean Absolute Percentage Error, or MAPE, is defined as the average of the percentage error between the actual and forecasted values. It is calculated as:

\[
\text{MAPE} = \frac{1}{n} \sum_{t=1}^{n} \left| \frac{A_t - F_t}{A_t} \right| \times 100
\]

where $A_t$ and $F_t$ are the actual and forecasted values at time $t$, respectively.

\subsubsection{Advantages of using MAPE}

\begin{itemize}
    \item It is easy to interpret since it is expressed as a    percentage.
    
    \item Furthermore, it allows for comparison across different time periods and scales
\end{itemize}

\subsubsection{Disadvantages of using MAPE}

The context of predicting the number of passengers that an
airline will carry poses some unique limitations to the usefulness of MAPE.

\begin{itemize}
    \item \textbf{Asymmetric Costs:} MAPE treats over-predictions and  under-predictions equally, which could be disastrous in this setting. Whereas overprediction may lead to some unnecessary operational costs, underprediction results in an inadequate fleet capacity, lost revenue and disgruntled customers. Therefore, our performance metric needs to take this into account.
    
    \item \textbf{Inability to Capture Peak Demand:} For fleet management and hiring, understanding peak demand is crucial. MAPE averages out the errors across all data points, which can mask important peaks or troughs in demand.
    
    \item \textbf{Instability with Low Values:} MAPE can become disproportionately large when actual values approach zero. Given that passenger numbers can vary significantly (from 150,000 to over 6,000,000), low passenger counts in certain months could lead to extremely high percentage errors, misleading decision-making.
\end{itemize}

\subsubsection{Alternatives to using MAPE}

Since the fleet requirements are constrained by the total number of passengers expected while human resources requirements are usually constrained by peak demand, there are some alternative evaluations to consider.

\begin{enumerate}
    \item \textbf{Mean Absolute Error (MAE):}
      \begin{itemize}
          \item
            MAE calculates the average of the absolute differences between the
            actual passenger numbers and the forecasted numbers.
      \end{itemize} 
\[
\text{MAE} = \frac{1}{n} \sum_{t=1}^{n} |A_t - F_t|
\]

\begin{itemize}
\item
  \textbf{Advantage:} It focuses on passenger count rather than
  percentages, thereby providing a direct assessment of the forecast's
  performance and making it easy for operational teams to understand and
  act upon.
\item
  \textbf{Disadvantage:} It still does not solve the problem of
  asymmetric consequences, since it reports the absolute value of the
  error. Also, it captures only the mean error and not the peak error.
\end{itemize}

\item
  \textbf{Weighted MAE:}

  \begin{itemize}
  \item
    We can introduce asymmetry in the calculation of MAE by taking a
    weighted mean of errors, assigning a relatively higher weight for
    underpredictions and a relatively lower weight for overpredictions.
  \item
    However, this loses the interpretability advantage that MAE offered,
    since we have now polluted the MAE with arbitrary weights.
  \end{itemize}
\item
  \textbf{Maximum Absolute Error (MaxAE):}

  \begin{itemize}
  \item
    MaxAE identifies the largest absolute error among all forecasted
    values compared to the actual values.
  \end{itemize}

\[
\text{MaxAE} = \max(|A_t - F_t|) \quad \text{for } t = 1, 2, \ldots, n
\]

\begin{itemize}
\item
  \textbf{Advantage:} It captures the worst-case scenario in terms of
  forecasting error. This is particularly critical for human resource
  planning.
\end{itemize}
\item \textbf{A Combination of MAE and MaxAE}:

  \begin{itemize}
      \item One can use MAE to plan fleet capacity based on total expected passengers and use MaxAE to ensure operational readiness for peak periods based on worst-case scenarios.
  \end{itemize}
  
\end{enumerate}

\newpage
\subsection{Part 3}

In order to determine whether or not the mean $\mu$ of the first-difference
series of \textit{`PASSENGERS CARRIED'} has changed in the pre-Covid and post-Covid era, we can utilize a statistical hypothesis test such as the \textbf{t-test for the difference in means} between two independent samples. 

\subsubsection{Hypotheses}

The hypotheses are as follows:

\begin{itemize}
\item
  \textbf{Null Hypothesis:} There is no difference in the means of the
  first differenced series of passenger numbers between the two periods.
    \[
    H_0: \mu_{\text{pre}} = \mu_{\text{post}}
    \]
\item
  \textbf{Alternative Hypothesis:} There is a difference in the means of
  the first differenced series of passenger numbers between the two
  periods.
    \[
    H_a: \mu_{\text{pre}} \neq \mu_{\text{post}}
    \]
\end{itemize}

\subsubsection{Test Statistic}

The test statistic for the independent samples t-test is calculated as follows:
\[
t = \frac{\bar{X}_{\text{pre}} - \bar{X}_{\text{post}}}{\sqrt{\frac{s_{\text{pre}}^2}{n_{\text{pre}}} + \frac{s_{\text{post}}^2}{n_{\text{post}}}}}
\]
Where:
\begin{itemize}
    \item \( \bar{X}_{\text{pre}} \) and \( \bar{X}_{\text{post}} \) are the sample means for pre-COVID and post-COVID, respectively.
    \item \( s_{\text{pre}}^2 \) and \( s_{\text{post}}^2 \) are the sample variances for the two groups.
    \item \( n_{\text{pre}} \) and \( n_{\text{post}} \) are the sample sizes for the two groups.
\end{itemize}

\subsubsection{Interpretation of Results}

The results of the t-test are interpreted as follows:

\begin{itemize}
\def\labelenumi{\roman{enumi}.}
\item
  The t-statistic is compared to a critical value from the
  t-distribution table based on our significance level (commonly $\alpha = 0.05$) and the degrees of freedom.
\item
 If $t$ is greater than the critical value, we reject the null hypothesis, suggesting a   significant difference in the mean number of passengers carried before and after the COVID period.
\item
  Rejecting the null hypothesis indicates a significant difference in
  the mean number of passengers carried between the two periods.
  Conversely, failing to reject the null hypothesis suggests no
  significant difference in passenger numbers before and after the
  pandemic exists.
\end{itemize}

\subsubsection{Assumptions}

The assumptions of the independent samples t-test are as follows:

\begin{itemize}
\item
  \textbf{Independence:} The observations in each group must be
  independent of one another.
\item
  \textbf{Normality:} The sample means must be approximately normally
  distributed. Since our first-difference series is weakly stationary,
  this assumption is reasonable under the Central Limit Theorem if the
  sample sizes are large enough.
\item
  \textbf{Homogeneity of Variances:} The variances of the two groups
  should be equal (or approximately equal). This can be tested using
  Levene's test for equality of variances.
\end{itemize}

\end{document}
